\section{Síntesis y ampliación}

Esta clase organizó los retrieval-augmented language models en un espectro continuo: desde frozen BM25/vector search, que inserta los documentos en el prompt; pasando por RePlug/reranker, que adaptan la retrieval a un frozen LM; luego RAG, REALM, Atlas y RETRO, que hacen llegar la task signal a más componentes; hasta generalizar la retrieval como active tool use, multimodal memory y systems-over-models.

\subsection{Marco unificado de comparación}

Cualquier sistema RAG puede describirse con una quíntupla: corpus/index $I$, retrieval policy $\pi$, retriever parameters $\theta_R$, fusion mechanism $F$ y generator parameters $\theta_G$. Al comparar artículos o productos, deben reportarse el estado de entrenamiento, el estado de despliegue y el costo de estos cinco elementos, en lugar de limitarse a decir «se usó RAG».

\begin{importantbox}{Diez conclusiones aplicables de esta clase}
\begin{enumerate}
  \item Evaluar por separado la hallucination, el generic error y la evidence contradiction;
  \item Reportar a la vez retrieval recall, evidence position, utilization y final correctness;
  \item Para una exact entity, conservar de preferencia la sparse signal; para la recuperación semántica, usar además dense/late interaction;
  \item Frozen RAG es un baseline sólido, pero hay que registrar el score/objective mismatch;
  \item La elección entre RePlug, reranker o joint training depende del generator access level;
  \item Reportar qué parámetros se actualizan, cuándo se refresca el index y si el corpus cambia;
  \item Controlar $k_r$, $k_c$ y la context position, y evitar ampliar el prompt a ciegas;
  \item La active retrieval debe optimizar el information value y el retrieval cost;
  \item El fine-tuning modifica la conducta y las representaciones; la retrieval modifica la per-request evidence; ambos pueden combinarse;
  \item Source policy, provenance, permissions y evaluation forman parte de la RAG architecture.
\end{enumerate}
\end{importantbox}

\subsection{Qué registrar al reproducir experimentos}

Como mínimo, registrar corpus snapshot, license/provenance, chunk size/overlap, embedding model, index type, retriever top-$k$, reranker, context packing, generator version, prompt, updated parameters, index refresh cadence, latency/cost y evaluation set. Sin esta información, es difícil reproducir los resultados aunque el nombre del modelo sea el mismo.

\begin{warningbox}{No presentar una instantánea de la clase de 2023 como un hecho de los productos de 2026}
Las valoraciones de esta clase sobre off-the-shelf retriever, vector database, hardware y tendencias de la industria se hicieron el 2023-12-05. Las notas conservan los mecanismos y los juicios históricos, pero no afirman con base en ellos las capacidades actuales de los productos, el panorama del mercado ni conclusiones jurídicas de 2026.
\end{warningbox}

\subsection{Lecturas adicionales}

Se recomienda leer en el orden de los mecanismos: DPR/ColBERT para entender la retrieval representation; RAG/FiD para entender la evidence marginalization y la fusion; kNN-LM/RETRO para entender cómo la non-parametric memory entra en el language modeling; REALM/Atlas para entender las aproximaciones del joint training; RePlug/In-Context RALM para entender el frozen generator; FLARE/Self-RAG/Toolformer para entender la active retrieval y el tool use; SILO/Lost-in-the-Middle para entender los límites de los datos y la evidence utilization.

\subsection{Resumen final}

El valor de la retrieval augmentation no se reduce a «darle más texto al modelo»: consiste en separar del modelo paramétrico único el estado del conocimiento, su actualización, sus fuentes y el cómputo. Un RAG de alta calidad debe alinear retriever, context, generator y source policy, y demostrar con una evaluation por capas que la evidencia correcta se encuentra, se coloca en una posición aprovechable y el generador realmente la usa, y que la salida final cumple con la tarea y con el ground truth.

\end{document}
